\chapter{Le secteur aurifère sénégalais et le projet Sabodala-Massawa}
\label{chap:senegal}

\section{Le cadre légal et fiscal du secteur minier sénégalais}

Le Sénégal dispose d'un potentiel minier diversifié --- phosphates, fer, zircon, attapulgite et or --- mais c'est essentiellement l'exploitation aurifère de la région de Kédougou, à la frontière avec le Mali et la Guinée, qui a transformé le secteur minier en source significative de recettes d'exportation et de recettes budgétaires depuis le début de la production commerciale en 2009. Le pays est membre de l'Initiative pour la Transparence dans les Industries Extractives depuis 2013 et publie, dans ce cadre, des rapports annuels de réconciliation des paiements du secteur extractif \citep{eiti_senegal}.

Le cadre légal applicable aux titres miniers a connu deux évolutions majeures~: le code minier de 2003, puis le code minier actuellement en vigueur, adopté par la loi n$^\text{o}$~2016-32 du 8 novembre 2016 \citep{senegal2016codeminier}. Ce dernier précise notamment les modalités d'octroi des titres de recherche et d'exploitation, le régime de la participation de l'État au capital des sociétés titulaires d'un permis d'exploitation --- fixée à 10\,\% à titre gratuit ---, ainsi que le cadre des conventions minières, qui peuvent stabiliser pour la durée du titre les dispositions fiscales et douanières applicables à un projet donné par rapport au droit commun.

C'est précisément ce mécanisme de convention minière qui régit le projet Sabodala, antérieur au code de 2016~: la convention minière de Sabodala, conclue initialement en 2005 et renégociée à plusieurs reprises (notamment en 2014 et en 2015), continue de s'appliquer au projet Sabodala-Massawa selon le principe de stabilité fiscale, tel que décrit à la section~\ref{sec:regime-sgo}. Cette coexistence entre le droit commun --- applicable aux nouveaux titres délivrés sous l'empire du code de 2016 --- et les conventions particulières antérieures constitue l'un des enjeux de transparence et de cohérence du cadre fiscal minier sénégalais, régulièrement souligné dans les rapports ITIE.

\section{Le projet Sabodala-Massawa}

\subsection{Présentation générale}

Le complexe Sabodala-Massawa est exploité par Sabodala Gold Operations S.A. (SGO), société de droit sénégalais détenue par Teranga Gold Corporation --- groupe canadien intégré en 2021 au sein d'Endeavour Mining --- aux côtés de la participation gratuite de 10\,\% de l'État sénégalais évoquée ci-dessus. Le site de Sabodala, en production depuis 2009, a été complété en 2016 par l'acquisition des actifs Massawa, auparavant détenus par Barrick Gold et Randgold Resources, ce qui a porté la superficie du permis d'exploitation et la base de ressources du projet à un niveau permettant d'envisager une exploitation intégrée sur une quinzaine d'années supplémentaires \citep{lycopodium2020pfs}.

L'étude de préfaisabilité (PFS) publiée en août 2020, qui constitue la principale source de données de ce rapport, présente un plan minier intégrant~:
\begin{itemize}
  \item l'exploitation à ciel ouvert de plusieurs gisements déjà en production (Sabodala, Golouma, Maki Medina, Goumbati West/Kobokoto) et de nouveaux gisements issus de l'acquisition Massawa (Sofia, Central Zone, North Zone, Delya, Masato, Niakafiri)~;
  \item l'exploitation souterraine de quatre gisements (Golouma West~1 et~2, Golouma South, Kerekounda)~;
  \item le traitement du minerai dans l'usine existante de Sabodala (procédé \emph{Whole-Ore-Leach}, WOL) complétée par une nouvelle unité de traitement du minerai réfractaire (\emph{Refractory Ore Treatment}, ROT) reposant sur la bio-oxydation (BIOX).
\end{itemize}

\subsection{Ressources et réserves}

Les réserves minérales prouvées et probables, estimées au 31 décembre 2019 sur la base d'un prix de l'or de 1\,250~\$US/once, s'élèvent à 75,8~millions de tonnes (Mt) à une teneur moyenne de 1,98~g/t, soit 4,82~millions d'onces (Moz) d'or contenu, réparties entre 49,0~Mt (1,84~Moz) pour le périmètre historique de Sabodala et 24,7~Mt (2,63~Moz) pour le périmètre Massawa \citep{lycopodium2020pfs}. Les ressources mesurées et indiquées, incluant les réserves, atteignent quant à elles 104,8~Mt à 2,05~g/t pour 6,9~Moz, auxquelles s'ajoutent 26,1~Mt de ressources inférées à 2,19~g/t pour 1,8~Moz.

\subsection{Calendrier de production}

Le PFS retient un horizon d'exploitation de 16,5~ans (mi-2020 à mi-2036), avec un traitement cumulé de 75,8~Mt de minerai (teneur moyenne 1,98~g/t) pour une production aurifère récupérée de 4,28~Moz, soit un taux de récupération moyen de 89\,\%. Le tableau~\ref{tab:profil-production} (chapitre~\ref{chap:methodo}) reproduit ce calendrier année par année~; il fait apparaître un profil de production relativement stable autour de 400~koz/an entre 2023 et 2026 --- période d'exploitation simultanée des unités WOL et ROT --- avant un repli progressif à partir de 2028 lorsque la contribution de l'usine ROT s'amenuise.

\subsection{Dépenses d'investissement et coûts d'exploitation}

Le PFS distingue les dépenses d'investissement de \textbf{développement} --- 409~millions \$US sur la durée de vie du projet, dont 257~millions \$US pour les phases~1 et~2 de mise à niveau de l'usine Massawa, 102~millions \$US pour le développement souterrain et 50~millions \$US pour la relocalisation de communautés --- et les dépenses de \textbf{maintien en exploitation} (\emph{sustaining capital}), évaluées à 241~millions \$US, principalement liées au renouvellement de la flotte minière et aux infrastructures de gestion des résidus.

Les coûts d'exploitation cumulés sur la durée de vie du projet sont estimés à 2\,560~millions \$US, répartis entre l'exploitation minière à ciel ouvert (1\,158~millions \$US), l'exploitation souterraine (155~millions \$US), le traitement WOL (709~millions \$US), le traitement ROT (234~millions \$US) et les coûts généraux et administratifs (304~millions \$US), auxquels s'ajoutent 34~millions \$US de coûts d'administration régionale. Le coût total de maintien en exploitation (\emph{All-In Sustaining Cost}, AISC) moyen sur les dix premières années (2021-2030) est estimé à 715~\$US/once \citep{lycopodium2020pfs}.

\subsection{Hypothèses économiques retenues}

Le PFS retient un prix de l'or de 1\,600~\$US/once pour l'analyse économique centrale --- sensiblement supérieur au prix de 1\,250~\$US/once utilisé pour la définition des teneurs de coupure des réserves --- ainsi qu'une hypothèse de prix de transport et de traitement/affinage net nul, l'or étant vendu directement sous forme de lingots. Sur cette base, le PFS estime la valeur actuelle nette après impôt et après prise en compte de la participation minoritaire de l'État à 2\,210~millions \$US (taux d'actualisation nominal de 0\,\%) et à 1\,622~millions \$US (taux d'actualisation de 5\,\%), pour une production cumulée vendue de 4\,285~koz sur la période 2020-2037.

\section{Le régime fiscal applicable au projet Sabodala-Massawa}
\label{sec:regime-sgo}

Le régime fiscal de SGO résulte de la combinaison du droit commun minier sénégalais et des dispositions particulières négociées dans le cadre de la convention minière de Sabodala et de ses avenants. Les principaux paramètres, tels que documentés dans le PFS \citep{lycopodium2020pfs}, sont les suivants~:

\begin{itemize}
  \item \textbf{Redevance minière}~: une redevance de 5\,\% est due sur la valeur des expéditions d'or. Ce taux résulte d'un accord conclu avec l'État en 2013, portant la redevance dite « \emph{net smelter royalty} » (NSR) initialement prévue par la convention de 2005 à ce niveau~;
  \item \textbf{Impôt sur les sociétés}~: le projet est assujetti au taux fédéral de droit commun de 25\,\%, l'exonération temporaire de huit ans accordée au démarrage de l'exploitation de Sabodala ayant expiré le 2~mai~2015. Au-delà de cette date, le régime fiscal de SGO est « stabilisé » par rapport au droit applicable à la date de signature de la convention initiale (23~mars~2005), sauf accord contraire des parties~;
  \item \textbf{Participation de l'État}~: conformément à l'accord d'actionnaires constitutif de SGO (novembre 2007), l'État sénégalais détient une participation gratuite de 10\,\% au capital de la société, les droits à dividende correspondants n'étant exigibles qu'après remboursement de l'investissement initial de l'actionnaire majoritaire et de la dette tierce~;
  \item \textbf{Retenue à la source sur les dividendes}~: appliquée au taux de droit commun lors de la distribution des bénéfices aux actionnaires non résidents~;
  \item \textbf{Droits de douane}~: SGO bénéficie, depuis le 16~novembre~2018, du statut d'Entreprise Franche d'Exportation (EFE), qui garantit une exonération des droits de douane --- y compris les prélèvements parafiscaux --- sur les équipements miniers, carburants, explosifs et produits chimiques importés, ainsi qu'une exonération de patente et de droits d'enregistrement, jusqu'au 15~octobre~2021. À titre de contrepartie, et indépendamment de cette exonération, SGO s'est engagée à contribuer annuellement à hauteur de 1,2~million \$US au développement des communautés riveraines~;
  \item \textbf{Taxe sur la valeur ajoutée}~: une exonération de collecte et de paiement de la TVA, accordée en février~2016 par avenant à la convention minière, est applicable jusqu'au 2~mai~2022~;
  \item \textbf{Fonds social}~: les onces produites à partir des gisements Massawa sont soumises à une contribution additionnelle de 0,5\,\% au profit d'un fonds de développement communautaire~;
  \item \textbf{Flux liés à l'acquisition de Massawa}~: un paiement initial de 15~millions \$US a été versé en 2020 au titre de la levée d'option de l'État sur une participation additionnelle dans Massawa, et un contrat d'approvisionnement en or (\emph{gold stream}) conclu en 2014 avec Franco-Nevada prévoit la livraison d'un volume fixe d'onces jusqu'en 2031, pour une valeur cumulée estimée à 140~millions \$US sur la durée du PFS.
\end{itemize}

Ces paramètres --- en particulier la redevance de 5\,\%, l'impôt sur les sociétés de 25\,\%, la participation gratuite de 10\,\% et l'exonération de droits de douane --- constituent le c\oe ur du scénario fiscal « Sénégal » construit dans le modèle utilisé pour ce rapport (chapitre~\ref{chap:methodo}). Les éléments plus spécifiques à la transaction Massawa (paiement initial, \emph{gold stream} Franco-Nevada, fonds social additionnel) ne sont, en revanche, pas représentés dans le modèle dans la mesure où ils relèvent d'arrangements contractuels propres à cette opération et non de paramètres généralisables du régime fiscal minier sénégalais~; leur ordre de grandeur est néanmoins rappelé ici pour mémoire et discuté au chapitre~\ref{chap:resultats}.
